\documentclass{article}
% \usepackage[czech]{babel}
\usepackage[T1]{fontenc}
\usepackage[utf8]{inputenc}
\usepackage{amsmath}
\usepackage{xcolor}

\renewcommand{\div}{\operatorname{div}}
\newcommand{\R}{\mathbb R}
\newcommand{\vc}[1]{\boldsymbol{#1}}
\def\d{{\rm{d}}}

\begin{document}

\title{Test č. 1}
\date{26. 11. 2026}
\author{}
\maketitle


\noindent

Napište slabou formulaci ustáleného problému vedení tepla na 1D oblasti $\Omega=[-1,1]$.
Tepelný tok je dán Fourierovým zákonem
\[
    q(x) = -\left(x^2\cos(x)+1\right)u'(x), \qquad \textrm{v }\Omega,
\]
a platí rovnice kontinuity
\[
    q'(x) = \sin(2x) \qquad \textrm{v }\Omega.
\]

\noindent
Na levé straně v bodě $-1$ je udržována stálá teplota $U_D$.\\
Na pravé straně v bodě $1$ je izolace o tepelné vodivosti $\sigma$, tepelný tok v tomto bodě je pak přímo úměrný rozdílu vnitřní
a vnější teploty, tedy $u(1)$, resp. $U_R$.

\begin{enumerate}
\item Zakreslete oblast $\Omega$, zaveďte souřadný systém, vyznačte okrajové podmínky.
\item Formulujte okrajové podmínky.
\item Definujte prostor, ve kterém leží slabé řešení $u$, a prostor testovacích funkcí.
\item Odvoďte slabou formulaci, aplikujte okrajové podmínky.
\item \textit{\textcolor{gray}{Definujte vhodné bilineární a lineární formy a zapište slabou formulaci pomocí nich.}}
\end{enumerate}


\vspace{1.5cm}
% \vfill
\emph{Bonus:}\\
Uvažujte slabou formulaci problému
\begin{align*}
    \sigma u(a)v(a) + \int\limits_{a}^{b} \lambda(x) & u'(x)v'(x) + u'(x)v(x) \,\d x = \\
        &\int\limits_{a}^{b} f(x)v(x) \,\d x - q_Nv(b) + \sigma u_R v(a) \qquad \forall v\in \mathrm{H}^1([a,b])
\end{align*}
Formulujte odpovídající silnou formulaci, tedy parc. diferenciální rovnici s okrajovými podmínkami.

\end{document}
